\documentclass[10pt,a4paper]{amsart}
\usepackage[margin=19mm]{geometry}
\usepackage{amsmath,amssymb,mathtools}
\usepackage[scr=rsfs,cal=euler]{mathalfa}
\usepackage{xfrac}
\usepackage[unicode,hidelinks,bookmarks=false]{hyperref}
\newcommand{\ie}{i.e.\ }
\newcommand{\cf}{cf.\ }
\newcommand{\resp}{respectively\ }
\newcommand{\Subsection}{\subsection}
\providecommand{\nobreakdash}{\nobreak\hbox{-}}
\setlength{\emergencystretch}{2em}
\multlinegap0pt
\usepackage{amssymb}
\usepackage{tikz}
\newtheorem{proposition}{Proposition}
\title{Multiparameter L\'evy white noise theory and applications}
\author{Olfa Draouil, Rahma Yasmina Moulay Hachemi, Bernt {\O}ksendal}
\date{Source: arXiv:2510.25525v2 (CC BY 4.0)}
\begin{document}
\maketitle

\noindent\emph{Source and licence.} Multiparameter L\'evy white noise theory and applications. Version arXiv:2510.25525v2 (CC BY 4.0), distributed by arXiv under the Creative Commons Attribution 4.0 International licence, http://creativecommons.org/licenses/by/4.0/, as recorded in the arXiv licence metadata for this exact version and shown on the abstract page https://arxiv.org/abs/2510.25525v2. Full licence notice: \texttt{licenses/arxiv-2510.25525-CC-BY-4.0.txt}.\par\smallskip

\noindent\textit{Editorial scope.} Counted unit: the article's proposition eq:K-bound (source0.tex lines 3106--3146). The article's complete original proof of it is delivered with it (source0.tex lines 3148--3613, one \texttt{proof} environment of 466 lines). No mathematics is written, repaired or re-derived: the statement and the proof are the article's own lines, in the article's own notation, and no retained line is substituted or re-flowed.  No further source-local context is needed: every cross-reference of every retained line resolves inside the two delivered spans.  Nothing was omitted from any retained span, so the package carries no omission record.  No retained line contains a citation, so the wrapper carries no bibliography and no bibliography entry.  No retained line contains a figure environment, an image-inclusion command or drawing code, so no image is delivered and the package declares no asset dependency.  Editorial formatting only, in addition: an AMS \texttt{amsart} preamble that restates the article's own declarations and macros used by the retained lines, namely the environments \texttt{proposition} on separate counters, together with the standard packages those lines need.  The retained environments print in this document as Proposition 1 in order of appearance, whereas the article prints the same items with the numbering of its own full document.  The complete native source of the article as distributed in the arXiv e-print for this exact version is bundled unchanged as \texttt{sources/188043/source0.tex}.
\begin{proposition}
Let $T,X>0$ and let $0<\alpha<1$. Assume that there exists an integer
$\ell_0\geq 0$ such that
\begin{equation}
\sup_{(s,z)\in[0,T]\times[0,X]}
\|\mathcal K(s,z)\|_{-1,-\ell_0}
\leq C_{\mathcal K}<\infty.
\label{eq:K-bound}
\end{equation}
Let $r$ be an integer satisfying
\begin{equation}
r>\ell_0+1.
\label{eq:r-condition}
\end{equation}
Assume moreover that
\[
(s,z)\longmapsto \mathcal K(s,z)
\]
is continuous from $[0,T]\times[0,X]$ into
$(\mathcal S)_{-1,-\ell_0}$.

Then the series
\begin{equation}
Y(t,x)=y_0+\sum_{n=1}^{\infty}Y_n(t,x)
\label{eq:Y-series}
\end{equation}
converges uniformly on $[0,T]\times[0,X]$ with values in
$(\mathcal S)_{-1,-r}$ and
\begin{equation}
Y\in
C\big([0,T]\times[0,X];(\mathcal S)_{-1,-r}\big).
\end{equation}
Furthermore, if $R_n$ denotes the remainder obtained after $n$ iterations
of the Volterra equation, then
\begin{equation}
\sup_{(t,x)\in[0,T]\times[0,X]}
\|R_n(t,x)\|_{-1,-r}
\longrightarrow0,
\qquad n\to\infty.
\end{equation}
\end{proposition}

\begin{proof}
We consider the Volterra equation
\begin{equation}
Y(t,x)
=
y_0+
\int_0^t\int_0^x
K(t,x;s,z)\diamond Y(s,z)\,dz\,ds,
\label{eq:Volterra}
\end{equation}
where
\begin{equation}
K(t,x;s,z)
=
\frac{(t-s)^{\alpha-1}}{\Gamma(\alpha)}
\mathcal K(s,z),
\qquad 0<\alpha<1.
\label{eq:K-definition}
\end{equation}

We first introduce the iterated kernels. For convenience, put
\[
u_0=(t,x),\qquad
u_j=(s_j,z_j),\quad j\geq1,
\]
and set $s_0=t$. We define
\begin{equation}
K(u_{j-1};u_j)
=
\frac{(s_{j-1}-s_j)^{\alpha-1}}
{\Gamma(\alpha)}
\mathcal K(s_j,z_j).
\label{eq:iterated-K}
\end{equation}

For $n\geq1$, define
\begin{equation}
\Delta_n(t,x)
=
\left\{
0<s_n<\cdots<s_1<t,\quad
0<z_n<\cdots<z_1<x
\right\}.
\label{eq:Delta}
\end{equation}

The $n$-th term of the iterated series is
\begin{equation}
Y_n(t,x)
=
y_0
\int_{\Delta_n(t,x)}
K(u_0;u_1)\diamond\cdots\diamond
K(u_{n-1};u_n)
\,du_n\cdots du_1.
\label{eq:Yn}
\end{equation}

Notice that the terminal point changes at every iteration. In particular,
the product in \eqref{eq:Yn} is
\[
K(t,x;s_1,z_1)
\diamond
K(s_1,z_1;s_2,z_2)
\diamond\cdots\diamond
K(s_{n-1},z_{n-1};s_n,z_n),
\]
and not a product in which $(t,x)$ is kept fixed in every factor.

We now estimate the iterated Wick products.

By \eqref{eq:K-bound} and \eqref{eq:iterated-K},
\begin{equation}
\begin{aligned}
\|K(u_{j-1};u_j)\|_{-1,-\ell_0}
&=
\frac{(s_{j-1}-s_j)^{\alpha-1}}
{\Gamma(\alpha)}
\|\mathcal K(s_j,z_j)\|_{-1,-\ell_0}
\\
&\leq
C_{\mathcal K}
\frac{(s_{j-1}-s_j)^{\alpha-1}}
{\Gamma(\alpha)}.
\end{aligned}
\label{eq:K-factor-estimate}
\end{equation}

Recall the Våge inequality: if
$F\in(\mathcal S)_{-1,-r}$ and
$G\in(\mathcal S)_{-1,-\ell_0}$ with
$r>\ell_0+1$, then
\begin{equation}
\|F\diamond G\|_{-1,-r}
\leq
A(r-\ell_0)
\|F\|_{-1,-r}
\|G\|_{-1,-\ell_0}.
\label{eq:Vage}
\end{equation}

Applying this inequality successively gives, for every $n\geq1$,
\begin{equation}
\begin{aligned}
&
\left\|
K(u_0;u_1)\diamond\cdots\diamond
K(u_{n-1};u_n)
\right\|_{-1,-r}
\\
&\qquad\leq
A(r-\ell_0)^{n-1}
\prod_{j=1}^{n}
\|K(u_{j-1};u_j)\|_{-1,-\ell_0}.
\end{aligned}
\label{eq:iterated-vage}
\end{equation}

Combining \eqref{eq:iterated-vage} with
\eqref{eq:K-factor-estimate}, we obtain
\begin{equation}
\begin{aligned}
&
\left\|
K(u_0;u_1)\diamond\cdots\diamond
K(u_{n-1};u_n)
\right\|_{-1,-r}
\\
&\qquad\leq
A(r-\ell_0)^{n-1}
C_{\mathcal K}^{\,n}
\prod_{j=1}^{n}
\frac{(s_{j-1}-s_j)^{\alpha-1}}
{\Gamma(\alpha)}.
\end{aligned}
\label{eq:wick-product-bound}
\end{equation}

Therefore, by the definition of $Y_n$ and the Bochner integral inequality,
\begin{equation}
\begin{aligned}
\|Y_n(t,x)\|_{-1,-r}
\leq{}&
|y_0|
A(r-\ell_0)^{n-1}
C_{\mathcal K}^{\,n}
\\
&\times
\int_{\Delta_n(t,x)}
\prod_{j=1}^{n}
\frac{(s_{j-1}-s_j)^{\alpha-1}}
{\Gamma(\alpha)}
\,du_n\cdots du_1.
\end{aligned}
\label{eq:Yn-bound1}
\end{equation}

The spatial and temporal integrations can be separated. First,
\begin{equation}
\int_{0<z_n<\cdots<z_1<x}
dz_n\cdots dz_1
=
\frac{x^n}{n!}.
\label{eq:spatial-simplex}
\end{equation}

On the other hand, by the iterated fractional integral identity,
\begin{equation}
\begin{aligned}
&
\int_{0<s_n<\cdots<s_1<t}
\prod_{j=1}^{n}
\frac{(s_{j-1}-s_j)^{\alpha-1}}
{\Gamma(\alpha)}
\,ds_n\cdots ds_1
\\
&\qquad=
\frac{t^{n\alpha}}
{\Gamma(n\alpha+1)}.
\end{aligned}
\label{eq:fractional-simplex}
\end{equation}

Consequently,
\begin{equation}
\boxed{
\|Y_n(t,x)\|_{-1,-r}
\leq
|y_0|
A(r-\ell_0)^{n-1}
C_{\mathcal K}^{\,n}
\frac{t^{n\alpha}}
{\Gamma(n\alpha+1)}
\frac{x^n}{n!}.
}
\label{eq:Yn-final-bound}
\end{equation}

Hence,
\begin{equation}
\begin{aligned}
\sup_{(t,x)\in[0,T]\times[0,X]}
\|Y_n(t,x)\|_{-1,-r}
\leq{}&
|y_0|
A(r-\ell_0)^{n-1}
C_{\mathcal K}^{\,n}
\\
&\times
\frac{T^{n\alpha}}
{\Gamma(n\alpha+1)}
\frac{X^n}{n!}.
\end{aligned}
\label{eq:uniform-Yn}
\end{equation}

We now prove the convergence of the numerical majorant. Since
$A(r-\ell_0)$, $C_{\mathcal K}$, $T$ and $X$ are fixed constants, it is enough to observe that
\[
\Gamma(n\alpha+1)n!
\]
grows faster than any exponential sequence. Consequently,
\begin{equation}
\sum_{n=1}^{\infty}
A(r-\ell_0)^{n-1}
C_{\mathcal K}^{\,n}
\frac{T^{n\alpha}X^n}
{\Gamma(n\alpha+1)n!}
<\infty.
\label{eq:majorant-convergence}
\end{equation}

It follows from \eqref{eq:uniform-Yn} and the Weierstrass
$M$-test that
\begin{equation}
\sum_{n=1}^{\infty}Y_n(t,x)
\end{equation}
converges uniformly on $[0,T]\times[0,X]$ in the norm of
$(\mathcal S)_{-1,-r}$.

It remains to prove continuity.

For every fixed $n$, the map
\[
(t,x)\longmapsto Y_n(t,x)
\]
is continuous with values in $(\mathcal S)_{-1,-r}$.

Indeed, by assumption,
\[
(s,z)\longmapsto\mathcal K(s,z)
\]
is continuous with values in $(\mathcal S)_{-1,-\ell_0}$.
The Wick product is continuous under the V$\mathring{a}$ge estimate, and the
integrand defining $Y_n$ is dominated by the integrable function
appearing on the right-hand side of \eqref{eq:wick-product-bound}.
Since $0<\alpha<1$, the singularity
$(s_{j-1}-s_j)^{\alpha-1}$ is integrable on the ordered simplex.
Therefore, the dominated convergence theorem for Bochner integrals
implies that
\[
Y_n\in
C\big([0,T]\times[0,X];(\mathcal S)_{-1,-r}\big).
\]

Since the series $\sum_{n\geq1}Y_n$ converges uniformly and each
$Y_n$ is continuous, its sum is continuous. Thus
\begin{equation}
\boxed{
Y\in
C\big([0,T]\times[0,X];(\mathcal S)_{-1,-r}\big).
}
\label{eq:Y-continuity}
\end{equation}

Because $[0,T]\times[0,X]$ is compact and $Y$ is continuous,
we immediately obtain
\begin{equation}
\boxed{
M_Y:=
\sup_{(t,x)\in[0,T]\times[0,X]}
\|Y(t,x)\|_{-1,-r}
<\infty.
}
\label{eq:Y-bound}
\end{equation}

We emphasize that \eqref{eq:Y-bound} is not an additional assumption.
It follows from the already established convergence of the series in
the space
$C([0,T]\times[0,X];(\mathcal S)_{-1,-r})$.

We now estimate the remainder. Iterating \eqref{eq:Volterra} $n$
times yields
\begin{equation}
Y(t,x)
=
y_0+\sum_{j=1}^{n}Y_j(t,x)+R_n(t,x),
\label{eq:remainder-decomposition}
\end{equation}
where
\begin{equation}
\begin{aligned}
R_n(t,x)
={}&
\int_{\Delta_{n+1}(t,x)}
K(u_0;u_1)\diamond\cdots\diamond
K(u_n;u_{n+1})
\\
&\qquad\diamond
Y(u_{n+1})
\,du_{n+1}\cdots du_1.
\end{aligned}
\label{eq:Rn-definition}
\end{equation}

There are $n+1$ kernel factors in \eqref{eq:Rn-definition}.
Since
\[
Y(u_{n+1})\in(\mathcal S)_{-1,-r}
\]
and
\[
\|Y(u_{n+1})\|_{-1,-r}\leq M_Y,
\]
successive applications of the V\ aage inequality yield
\begin{equation}
\begin{aligned}
&
\left\|
K(u_0;u_1)\diamond\cdots\diamond
K(u_n;u_{n+1})
\diamond Y(u_{n+1})
\right\|_{-1,-r}
\\
&\qquad\leq
M_Y
A(r-\ell_0)^{n+1}
C_{\mathcal K}^{\,n+1}
\prod_{j=1}^{n+1}
\frac{(s_{j-1}-s_j)^{\alpha-1}}
{\Gamma(\alpha)}.
\end{aligned}
\label{eq:Rn-integrand-bound}
\end{equation}

Consequently,
\begin{equation}
\begin{aligned}
\|R_n(t,x)\|_{-1,-r}
\leq{}&
M_Y
A(r-\ell_0)^{n+1}
C_{\mathcal K}^{\,n+1}
\\
&\times
\int_{\Delta_{n+1}(t,x)}
\prod_{j=1}^{n+1}
\frac{(s_{j-1}-s_j)^{\alpha-1}}
{\Gamma(\alpha)}
\,du_{n+1}\cdots du_1.
\end{aligned}
\label{eq:Rn-bound1}
\end{equation}

Using again the identities
\begin{equation}
\int_{0<z_{n+1}<\cdots<z_1<x}
dz_{n+1}\cdots dz_1
=
\frac{x^{n+1}}{(n+1)!}
\end{equation}
and
\begin{equation}
\begin{aligned}
&
\int_{0<s_{n+1}<\cdots<s_1<t}
\prod_{j=1}^{n+1}
\frac{(s_{j-1}-s_j)^{\alpha-1}}
{\Gamma(\alpha)}
\,ds_{n+1}\cdots ds_1
\\
&\qquad=
\frac{t^{(n+1)\alpha}}
{\Gamma((n+1)\alpha+1)},
\end{aligned}
\end{equation}
we obtain
\begin{equation}
\boxed{
\begin{aligned}
\|R_n(t,x)\|_{-1,-r}
\leq{}&
M_Y
A(r-\ell_0)^{n+1}
C_{\mathcal K}^{\,n+1}
\\
&\times
\frac{t^{(n+1)\alpha}}
{\Gamma((n+1)\alpha+1)}
\frac{x^{n+1}}
{(n+1)!}.
\end{aligned}
}
\label{eq:Rn-final}
\end{equation}

Therefore,
\begin{equation}
\begin{aligned}
\sup_{(t,x)\in[0,T]\times[0,X]}
\|R_n(t,x)\|_{-1,-r}
\leq{}&
M_Y
A(r-\ell_0)^{n+1}
C_{\mathcal K}^{\,n+1}
\\
&\times
\frac{T^{(n+1)\alpha}}
{\Gamma((n+1)\alpha+1)}
\frac{X^{n+1}}
{(n+1)!}.
\end{aligned}
\label{eq:Rn-uniform}
\end{equation}

The right-hand side of \eqref{eq:Rn-uniform} converges to zero as
$n\to\infty$. Hence
\begin{equation}
\boxed{
\sup_{(t,x)\in[0,T]\times[0,X]}
\|R_n(t,x)\|_{-1,-r}
\longrightarrow0.
}
\label{eq:Rn-to-zero}
\end{equation}

Thus the infinite iterated expansion is valid in
$(\mathcal S)_{-1,-r}$ and satisfies the original Volterra equation.

In particular,
\[
Y(t,x)
=
y_0+\sum_{n=1}^{\infty}Y_n(t,x)
\]
converges uniformly on $[0,T]\times[0,X]$, and
\[
Y\in
C\big([0,T]\times[0,X];(\mathcal S)_{-1,-r}\big).
\]

Finally, we stress that no smallness condition of the form
\[
A(r-\ell_0)C_{\mathcal K}<1
\]
is required. The convergence follows from the Gamma-function and
factorial terms
\[
\frac{T^{n\alpha}}{\Gamma(n\alpha+1)}
\qquad\text{and}\qquad
\frac{X^n}{n!},
\]
which dominate the exponential growth produced by the successive
applications of the V$\mathring{a}$ge inequality.
\end{proof}

\end{document}
